\section{Experimental Methodology}\label{sec:method}

We fixed the protocol described below, including the hypotheses, the refutation criteria, and all parameters, in writing on 6 November 2026, before running experiments E1 to E3. The MILP scaling study had been run earlier, on 2 November 2026, and we report it only as classical context. Table~\ref{tab:parameters} summarizes the parameters of every experiment.

\begin{table}[t]
\centering
\caption{Parameters of the experiments.}\label{tab:parameters}
\footnotesize
\begin{tabular}{@{}l>{\raggedright\arraybackslash}p{0.74\linewidth}@{}}
\toprule
Experiment & Setting \\
\midrule
E1 (H1) & Layered family, unit costs, $w\in\{2,\dots,5\}$, $k\in\{1,\dots,w-1\}$, seeds 0--2 (30 instances, $N$ from 9 to 974{,}121); $R=200$ runs per method and instance \\
E2 & Random family, $\rho=2$, unit costs, $n\in\{3,\dots,7\}$, $k\in\{1,2\}$ with $k<\kappa$, seed 0 (9 instances, $m$ from 3 to 14); transpilation to \texttt{FakeFez} at optimization level 1 \\
E3 (H2) & Random family, $n=m=3$, unit costs, $k=1$, $M=1$; $j\in\{0,1,2,3\}$ Grover iterations; 4{,}096 shots per circuit; optimization level 3; ideal simulator, \texttt{FakeFez}, and \texttt{ibm\_fez} \\
Feasible set & Count of Section~\ref{sec:interface-space}; random family, $\rho=2.5$, $n\in\{5,6,7,8,10\}$, seeds 0--2, costs in $\{1,\dots,5\}$, $B$ equal to 50\% of the cost of the cheapest $s$--$t$ cut \\
MILP study & Random family with $3n$ arcs, $n\in\{10,20,40,\dots,5{,}120\}$, and layered family with $w\in\{4,6,8,12,16,24,32,48,64\}$; costs in $\{1,\dots,5\}$; $B$ equal to 25\%, 50\%, or 75\% of the cost of the cheapest $s$--$t$ cut; seeds 0--2; 20\,s per instance \\
All & Capacities in $\{1,\dots,10\}$; $Q=22.5\sqrt{N}+1.4\log_2^2N$; $\lambda=6/5$; MILP reference limit 600\,s; 2{,}000 bootstrap resamples; tests at level 0.05 \\
\bottomrule
\end{tabular}
\end{table}


\subsection{Instances}\label{sec:method-instances}

The instances come from two seeded generators, and each instance is identified by a SHA-256 hash of its arcs, capacities, costs, and budget. The \emph{layered} family of width $w$ has a source, $w$ layers of $w$ nodes each, and a sink, so $n=w^2+2$. The source connects to every node of the first layer, every node of the last layer connects to the sink, and every other node connects to $\min(3,w)$ distinct nodes of the next layer chosen uniformly at random. The \emph{random} family with $n$ nodes contains a random Hamiltonian $s$--$t$ path, which guarantees connectivity, and random arcs up to $\lfloor\rho n\rfloor$ arcs in total. Capacities are integers drawn uniformly from $\{1,\dots,10\}$. Every instance used to test the interface has unit costs and budget $k<\kappa$, where $\kappa$ is the smallest number of arcs that disconnects $s$ from $t$, so that $\Fstar>0$ (Section~\ref{sec:bg-mfnip}). We also checked that $\Fstar<F(0)$, so that interdiction makes a difference. Figure~\ref{fig:example} shows an example.

\begin{figure}[t]
\centering
\begin{tikzpicture}[
  font=\footnotesize,
  vertex/.style={circle, draw=black, line width=0.4pt, fill=white, inner sep=0pt, minimum size=4.4mm},
  arc/.style={draw=black, line width=0.45pt, -{Stealth[length=1.6mm, width=1.1mm]}},
  removed/.style={arc, draw=qaccent, line width=0.8pt, dash pattern=on 3pt off 1.5pt},
  capacity/.style={fill=white, inner sep=1pt},
]
\node[vertex] (s) at (0.840,1.300) {$s$};
\node[vertex] (v1) at (2.400,2.600) {1};
\node[vertex] (v2) at (2.400,1.300) {2};
\node[vertex] (v3) at (2.400,0.000) {3};
\node[vertex] (v4) at (4.800,2.600) {4};
\node[vertex] (v5) at (4.800,1.300) {5};
\node[vertex] (v6) at (4.800,0.000) {6};
\node[vertex] (v7) at (7.200,2.600) {7};
\node[vertex] (v8) at (7.200,1.300) {8};
\node[vertex] (v9) at (7.200,0.000) {9};
\node[vertex] (t) at (8.760,1.300) {$t$};
\draw[arc] (s) -- (v1);
\draw[removed] (s) -- (v2);
\draw[removed] (s) -- (v3);
\draw[arc] (v1) -- (v4);
\draw[arc] (v1) -- (v5);
\draw[arc] (v1) -- (v6);
\draw[arc] (v2) -- (v4);
\draw[arc] (v2) -- (v5);
\draw[arc] (v2) -- (v6);
\draw[arc] (v3) -- (v4);
\draw[arc] (v3) -- (v5);
\draw[arc] (v3) -- (v6);
\draw[arc] (v4) -- (v7);
\draw[arc] (v4) -- (v8);
\draw[arc] (v4) -- (v9);
\draw[arc] (v5) -- (v7);
\draw[arc] (v5) -- (v8);
\draw[arc] (v5) -- (v9);
\draw[arc] (v6) -- (v7);
\draw[arc] (v6) -- (v8);
\draw[arc] (v6) -- (v9);
\draw[arc] (v7) -- (t);
\draw[arc] (v8) -- (t);
\draw[arc] (v9) -- (t);
\node[capacity] at (1.620,1.950) {4};
\node[capacity, text=qaccent] at (1.620,1.300) {9};
\node[capacity, text=qaccent] at (1.620,0.650) {6};
\node[capacity] at (3.000,2.600) {1};
\node[capacity] at (3.000,2.275) {8};
\node[capacity] at (3.000,1.950) {8};
\node[capacity] at (3.000,1.625) {9};
\node[capacity] at (3.000,1.300) {2};
\node[capacity] at (3.000,0.975) {1};
\node[capacity] at (3.000,0.650) {9};
\node[capacity] at (3.000,0.325) {1};
\node[capacity] at (3.000,0.000) {6};
\node[capacity] at (5.400,2.600) {1};
\node[capacity] at (5.400,2.275) {3};
\node[capacity] at (5.400,1.950) {5};
\node[capacity] at (5.400,1.625) {5};
\node[capacity] at (5.400,1.300) {5};
\node[capacity] at (5.400,0.975) {1};
\node[capacity] at (5.400,0.650) {1};
\node[capacity] at (5.400,0.325) {2};
\node[capacity] at (5.400,0.000) {1};
\node[capacity] at (7.980,1.950) {7};
\node[capacity] at (7.980,1.300) {6};
\node[capacity] at (7.980,0.650) {7};
\end{tikzpicture}

\caption{Layered instance with $w=3$ ($n=11$, $m=24$, unit costs, $k=2$, $N=301$) and arc capacities. The dashed arcs form the only optimal plan, which reduces the maximum flow from $F(0)=19$ to $\Fstar=4$.}\label{fig:example}
\end{figure}

Experiment E1 measures how the cost to the optimum scales on 30 layered instances (H1). Experiment E2 counts the resources of the interface on nine random instances with 3 to 14 arcs. Experiment E3 runs amplitude amplification on hardware (H2). Its instance is the first seed of the random family with $n=3$ and $m=3$ whose instance has $k=1<\kappa$, a single optimal plan ($M=1$), and $0<\Fstar<F(0)$, which we selected automatically. Both E3 variants use the same truth-table oracle, which marks only the optimal plan (threshold $\tau=\Fstar+1$). The \emph{interface} variant uses the Dicke preparation over $\mathcal{F}_1$ ($N=4$). The \emph{hypercube} variant uses $H^{\otimes 3}$ ($N=8$) and enforces the budget only implicitly through the oracle. We accessed IBM processors through the Open Plan of the IBM Quantum Platform, which offered a monthly quota of 10 minutes of processor time when this paper was written. The eight circuits of E3 formed a single job that used a few seconds of this quota. The MILP scaling study solved random and layered instances with general costs to find the largest networks for which the MILP certifies the optimum within 20 seconds.

\subsection{Methods}\label{sec:method-methods}

\paragraph{Minimum finding at the query level.} The quantum algorithm is Algorithm~\ref{alg:dh}. Simulating it does not require simulating the oracle gate by gate. For each instance, we enumerate $\Fk$ once and compute $F(x)$ for every plan with a classical maximum flow routine, and this table defines the oracle. This is valid because we measure the number of queries and evaluations rather than the time of each one. Given a threshold $\tau$ and $j$ iterations, the outcome is drawn from the exact distribution of amplitude amplification from $\Ak$. With probability $\sin^2((2j+1)\theta)$, where $\sin^2\theta=M_\tau/N$, the outcome is a marked plan chosen uniformly among the marked ones, and otherwise it is an unmarked plan chosen uniformly (Section~\ref{sec:bg-grover}). This \emph{reduced} simulator draws each outcome in $\bigO(\log N)$ time after sorting the flows, which makes $N\approx 10^6$ reachable. We validated it against a \emph{state-vector} simulator that stores the $N$ amplitudes and applies the oracle and the reflection explicitly (Section~\ref{sec:method-validation}). In E1, each run stops at the first measurement of an optimal plan, and we record the queries and evaluations up to that point and whether their sum fits within the budget $Q$ of Proposition~\ref{prop:advantage}. Because the random choices of a run do not depend on when the run is interrupted, this record is exactly the success indicator of a run interrupted at $Q$.

\paragraph{Classical rivals.} Exhaustive enumeration evaluates every plan of $\Fk$, first by cardinality and then in lexicographic order, and certifies the optimum with exactly $N$ evaluations. Random-order search evaluates the plans of $\Fk$ in uniformly random order until it finds an optimal plan. It has the same kind of probabilistic guarantee as the quantum algorithm, and its expected cost is $(N+1)/(M+1)$ evaluations. Both rivals search the same set $\Fk$ as the quantum algorithm.

\paragraph{Reference optimum.} For each instance, we solve the MILP~\eqref{eq:milp} with HiGHS~\cite{highs2018} through SciPy~\cite{scipy2020}, which certifies $\Fstar$ independently of the enumeration. Every solution returned by any method is re-evaluated by an independent maximum flow computation before it is counted.

\paragraph{Circuits.} We implemented the preparation $\Ak$, the truth-table oracle, the arithmetic oracle, and the Grover iteration (built with \texttt{grover\_operator}) in Qiskit~\cite{qiskit2024}. Logical resources are counted after transpilation to the basis $\{\texttt{cx},\texttt{u}\}$ without connectivity constraints, and physical resources after transpilation to \texttt{FakeFez}, the noise model of the 156-qubit IBM Heron processor \texttt{ibm\_fez}. The circuits run with the \texttt{SamplerV2} primitive of Qiskit IBM Runtime on the ideal Aer simulator, on \texttt{FakeFez}, and on the IBM processor.

\subsection{Metrics}\label{sec:method-metrics}
\begin{itemize}
\item The \emph{cost to the optimum} is the number of queries plus evaluations until an optimal plan is first measured (minimum finding), the number of evaluations until an optimal plan is found (random-order search), or $N$ (enumeration). For each instance, we use the mean over the $R$ runs.
\item The \emph{scaling exponent} $\beta$ is the least-squares slope of $\log_{10}$ of the mean cost against $\log_{10}N$ over the 30 instances. We also fit it against $\log_{10}(N/M)$ to separate the effect of the number of optimal plans, and we obtain 95\% confidence intervals from 2{,}000 bootstrap resamples of the instances. We fit no exponent for enumeration, whose cost is $N$ by definition.
\item \emph{Success within the theoretical budget} is the fraction of runs whose cost to the optimum is at most $Q$, which Proposition~\ref{prop:advantage} guarantees to be at least $1/2$.
\item \emph{Resources} are the number of qubits, the number of two-qubit gates, and the circuit depth, in logical and transpiled form.
\item \emph{Amplification} is the frequency of the optimal plan among the shots. We compare it with the uniform baseline $1/N$ using a one-sided binomial test, with the theoretical probability $\sin^2((2j+1)\theta)$ using the ratio of the two and a 95\% Wilson interval, and with the frequency on \texttt{FakeFez} using a two-sided Fisher exact test. For the hardware job, we also record the processor time reported by IBM and the wall-clock time from submission to result, which includes the queue.
\item The \emph{classical context} consists of the wall-clock time of the MILP and of single-threaded enumeration of $\Fk$ on the instances of E1, and of the largest network size up to which the MILP certified all three seeds of every size in the MILP scaling study.
\end{itemize}

\subsection{Hypotheses and refutation criteria}\label{sec:method-hypotheses}
\paragraph{H1 (scaling in simulation).} H1 is refuted if the fitted exponent $\beta$ of minimum finding falls outside $[0.4,0.6]$. The decision uses the point estimate, and we report the confidence interval. Because the reduced simulator draws each outcome from the exact amplification probability, an exponent close to $0.5$ follows from the theory under test whenever the control loop and the cost accounting are implemented correctly. H1 is therefore a correctness check of the implementation, not a test of a physical oracle. The registered interval is also wide for a theoretical value of $0.5$, so we additionally report whether the confidence interval contains $0.5$, a check that was not part of the registered criterion.

\paragraph{H2 (hardware).} On an IBM processor, after the optimal number of iterations of each variant ($j=1$ for the interface variant, where theory gives probability one, and $j=2$ for the hypercube variant, where it gives $0.945$), the frequency of the optimal plan is significantly greater than $1/N$ (one-sided binomial test at level $0.05$). H2 is refuted otherwise. This registered baseline is weak, because any circuit that works at all should pass it. We therefore also compare the hardware frequency with the theoretical probability and with the frequency on \texttt{FakeFez}. We defined these comparisons after the hardware run, so we report them as exploratory and not as tests of a hypothesis fixed in advance.

\subsection{Validation before measurement}\label{sec:method-validation}
No number of E1 was produced before the following checks, implemented as automated tests in the repository, had passed.
\begin{enumerate}
\item The reduced and state-vector simulators give the same outcome distribution, up to $10^{-10}$, for every threshold and for $0$ to $2\sqrt{N}+1$ iterations, with $N\le 2^{10}$.
\item The Qiskit preparation $\Ak$ produces the uniform superposition over $\Fk$ for all $m\le 6$ and $0\le k\le m$, and the Qiskit Grover circuit reproduces the distribution of the query-level model up to $10^{-10}$ for $0$ to $3$ iterations, which links the circuits to the simulator used in E1.
\item The arithmetic oracle flips the sign of exactly the pairs $(x,\alpha)$ with $C(\alpha,x)<\tau$ on all $2^{7}$ basis states of an instance with five arcs and returns every ancilla to zero.
\end{enumerate}
The last check exposed an error in the first version of our arithmetic oracle, in which the circuit synthesizer treated input qubits that had not yet been acted on as clean ancillas. We rebuilt the oracle with an explicit work register before counting any resources. For the hardware experiment, each circuit first ran on the ideal simulator and was accepted only if every frequency of the optimum lay within three standard errors of theory. It then ran on the noisy simulator and only afterwards on the processor.

\subsection{Computational environment}
All classical computations ran on a laptop with an Intel Core i5-1335U processor and 14.8\,GB of memory, using Python 3.14.4, Qiskit 2.6.1, Qiskit Aer 0.18.2, and Qiskit IBM Runtime 0.51.0. Wall-clock times were measured on a single thread, except for the MILP, for which HiGHS chooses the number of threads.

\subsection{Code and data availability}\label{sec:method-code}
The code, the experiment records, and the notebooks that produce every figure and table are publicly available at \url{https://github.com/VinTesta/q-mfnip}. The code is a Python package with a locked environment (\texttt{uv.lock}), and every number in Section~\ref{sec:results} is read from files written by the notebooks. The repository contains the package (\texttt{src/qmfnip/}), the validation tests of Section~\ref{sec:method-validation} (\texttt{tests/}), one notebook per experiment (\texttt{notebooks/}), and one output folder per experiment with its manifest, environment, instances, run records, summaries, and the identifier of the hardware job (\texttt{outputs/}). It also contains the experimental protocol, the list of reported results with the file each number comes from, and the queries of the literature search (\texttt{docs/}).
